\section{Introduction}
\label{sec:intro}

Computer-aided design (CAD) is the foundational representation in engineering and manufacturing.
Unlike meshes or point clouds, CAD programs are editable, interoperable with simulation and manufacturing tools, and capture the design intent behind a part rather than just its surface geometry.
This has motivated a growing line of research on parametric CAD generation~\cite{deepcad,skexgen,cadsignet} and, more recently, code-oriented approaches that treat CAD as a structured executable language amenable to language model training~\cite{cadrille,cadrecode}.

Despite this progress, complex CAD generation remains fundamentally difficult.
A real industrial part is the product of dozens of ordered sketch-and-extrude operations, each constrained by plane orientation, arc parameterization, and feature adjacency.
A small local mistake---a sign flip in a coordinate, an arc midpoint that is a few millimeters off its intended circle---propagates through the construction history and produces a wrong geometry that may look superficially plausible but differs materially from the target.
Direct generation of such geometry is hard to validate: a shape can pass visual inspection while remaining structurally incorrect or unusable as training supervision.

We argue that the core difficulty is not only \emph{how to generate} CAD programs, but \emph{how to know when a generated program is correct}.
Our central observation is that CAD is unusually well-suited for a precise reformulation: CAD artifacts can be represented as deterministically executable programs, and executable programs admit strong automated validation signals that require no human annotation.
Rather than treating CAD generation as open-ended geometry synthesis, we cast it as \textbf{verifiable program synthesis} (Figure~\ref{fig:teaser}).
A language model generates a CadQuery program; the program is executed in a sandboxed subprocess into a STEP file; the STEP file is validated against the ground-truth geometry via 3D IoU computed by OCCT boolean operations.
This converts an otherwise weakly supervised generative problem into one with explicit, computable acceptance criteria.

A second challenge is systems design at scale.
In practice, a complex engineering dataset is assembled over many batch runs, across multiple providers, with a heterogeneous failure mode distribution.
Some parts fail because the LLM misidentifies the sketch plane; others because \texttt{radiusArc} in CadQuery is ambiguous for non-quarter arcs; others because the ground-truth STEP records only the last extrude body rather than all accumulated bodies.
No single repair heuristic addresses all of these.
We therefore organize the pipeline as a collection of reusable \textbf{skills} (Figure~\ref{fig:skills}) rather than a monolithic workflow: separate skills for generation, execution, verification, repair, manual checking, and dataset assembly can be composed, replaced, or extended independently without redesigning the overall system.

\paragraph{Contributions.}
\begin{enumerate}[leftmargin=*,topsep=2pt,itemsep=1pt]
  \item We reformulate complex CAD generation as verifiable program synthesis and identify executable CAD code as the representation that makes strong automated verification tractable.
  \item We introduce a verification-guided curation pipeline (Figure~\ref{fig:pipeline}) that executes generated CAD programs and filters them through geometric, visual, and metadata checks.
  \item We propose a skill-based modular system that handles heterogeneous failure modes and recovers hard-tail cases through targeted repair and manual checking.
  \item We demonstrate the system builds a corpus of 1,646 verified program--geometry pairs with IoU $\geq 0.99$, covering complex multi-operation industrial parts that are beyond the reach of fully automated synthesis alone.
\end{enumerate}

\begin{figure*}[t]
  \centering
  \includegraphics[width=\linewidth]{fig1_teaser.png}
  \caption{\textbf{From Complex CAD Generation to Verifiable Data Generation.}
  (\textit{Left}) Complex CAD construction is long-horizon and error-prone: a sign flip in a sketch coordinate or an arc midpoint computed from the wrong formula produces a shape that looks plausible but differs materially from the target. (\textit{Center}) Our reformulation: generate an executable CadQuery program rather than direct geometry. (\textit{Right}) Execute the program, compare to GT STEP via 3D IoU, and accept if IoU $\geq 0.99$. This converts an annotation-free verification signal into a precise acceptance gate, producing high-confidence training-ready program--geometry pairs.}
  \label{fig:teaser}
\end{figure*}
